\chapter{Perhaps hopelessness is (The Only Thing I Have Left to Guide Me)}
We begin at the end, with the wheel of fate \cite{ref78}. It is everywhere in this episode. The great grinding gears of Walpurgisnacht are the cycles of hope and despair that transform the magical girls (back) into witches, and they are the endless cycles of reset time Homura creates. Both are the wheel of fate, as is the round clockwork buckler Homura uses to travel in time.

At the far end of the episode from them, Kyubey explains that Homura created Madoka. With every rebirth, Madoka carries the karmic burdens of the previous cycle into the next \cite{ref79}. In other words, though Madoka carries no concrete memories from time to time, only a few vague impressions no doubt arising from her deep connections to the one person who does remember previous timelines, and though she is physically recreated in each timeline, she still maintains the connections to the world of the previous timeline. And since these worlds are being created and destroyed for her, it is the weight of the entire world she carries. 

And it is a weight. In the end, these connections to a world hurtling headlong to destruction can lead only to suffering \cite{ref80}. The rising entropy of the world feeds into Madoka, transforming her in successive timelines from an inexperienced, but outgoing and confident, magical girl to an uncertain and unpowered sidekick whose associated heroines keep dying. 

And why does this happen? Because Homura tried to protect her. Homura is a Christian (or, at least, attended a Christian school). Like Kyoko, she has absorbed the belief that it is possible for one person to save another—that there is something to be saved from, and somewhere to be saved to. It is a fundamentally dualistic proposition—here is bad, but there is good. But as Kyubey has made clear with his talk of karmic balance, we are in a Buddhist world. What both Kyubey and Homura have failed to understand is that in a Buddhist universe, the decay of the material world is an illusion because all distinctions are illusions \cite{ref81}. The past is the present is the future. Decay is life. Magical girls are witches. The material is immaterial, and the other is the self. Homura cannot save Madoka because there is no Madoka to save and nothing from which to save her. All things are One—and since this is a fictional story, that One is the narrative itself, and in turn the gestalt entity within which that narrative takes form, the implied author.

We no longer have a Mami to contrast with Homura, but we do have the woman who kept being paralleled to Mami in the early episodes, Junko, who has two very important scenes in this episode. In the first, she discusses Sayaka's disappearance and its effect on Madoka with Madoka's English teacher. In this scene, she seems to parallel Sayaka more than Mami: she is gently chastened for her inability to stand by and do nothing, for insisting on doing rather than being, and a shot from the teacher's point of view is bathed in blue light and framed to focus on her hair clip, with the result that she even looks like Sayaka. 

Hovering above Junko throughout the conversation with Madoka's teacher is a reproduction of Michelangelo's The Creation of Adam. Notably, the position of the two characters aligns Junko with God and the teacher with Adam; at the same time, the red lighting on the teacher's side and blue lighting on Junko's causes God's red mantle (which has been compared by scholars to both a uterus \cite{ref82} and a brain \cite{ref83}) to be barely visible; it is instead Adam who appears wrapped safely in red warmth. This image foreshadows Junko's second scene in the episode, in the shelter, when she realizes that she has to stop protecting Madoka and start trusting her. It is a tragic scene, culminating in possibly the single densest shot in the series, a simple image of Junko's mom-jeans-and-sweater-clad crotch, lower abdomen, and hand, Madoka's mother reduced to a womb. The hand reaches out as if to grab Madoka, pull her inside—and then stops, and we cut back to Junko's face. The woman will not be idealized (read: reduced and objectified) as a mother; she overrides her own mothering tendency and makes the active choice not to act, to allow Madoka to risk the life Madoka owns. Watching, we feel Junko's pain as she lets go, but we also feel the profound respect and trust she is showing toward her daughter.

But Homura has been doing the opposite, pursuing a Christian ideal of salvation, in which the higher protects the lower, a permanent womb. She is trying to block Madoka from choosing the self-sacrifice she knows Madoka would make; if Junko is showing respect and trust toward Madoka, what is Homura showing toward her?

The answer, as of all people the supposedly unempathic and unemotional Kyubey notes (revealing once again that he has plenty of intellectual empathy, just no emotional empathy, a not entirely inaccurate first-order description of a sociopath \cite{ref84}), is that it has long ago ceased to be about Madoka. Homura herself described her journey to protect Madoka as a labyrinth, and the inside of her apartment resembles nothing so much as the surreal interior of a witch's labyrinth. Extradiegetically, magical girls began as witches; diegetically, they all end as witches. Homura can transcend time. In Homura, all things are one: the past, present, and future become a single wheel. She is, in other words, both magical girl and witch, trapped in a temporal labyrinth that is the narrative of the series itself. In that role as magical girl/witch, she brings both wishes and curses, at once protecting Madoka and ensuring her destruction. 

But what transforms a magical girl into a witch is the transition from hope to despair, from the optimism that things will work out to the realization that death and decay are inevitable and inescapable, and Homura has never experienced that transition. Homura has never had hope. She was a weak and sickly child who became a magical girl out of desperation and a sense of duty, wishing to become not Madoka's savior but her protector. Even in her wishes, she does not imagine that she will actually succeed, only that she will try. She is entirely and innately hopeless, and as such she is immune to despair—because like all other binaries, hope and despair are one. Thus it is not hopelessness that finally breaks Homura; rather, the realization that she is making things worse causes her to doubt her path for the very first time. For her, it is the transition from determination to doubt that threatens to bring about her transformation.

And there laughing at her is Walpurgisnacht, the witches' Sabbath \cite{ref85}, the wheel of fate, the Harlequin. A theatrical figure, the Harlequin fulfills the promise of the rising curtain at the beginning of the first episode—and, too, the theater plays the same show every night, with only slight variations, reflecting the wheel of fate and Homura's endless looping. In the harlequinades of the old commedia dell'arte, the Harlequin is a trickster figure, mocking all authority and order, and in particular making a mockery of tragic tales of doomed romance \cite{ref86}. Like the Harlequin, Walpurgisnacht will dance forever with Homura, laughing at her and snatching away Madoka, her love, defying rules or containment, a living symbol of the irreducible unpredictability and chaos of life. There is no higher and lower in the face of a trickster, only people. Rules are broken. Systems come crashing down.

And it is in this moment that Madoka arrives to make her wish.
